% ECONOMICS_PROBLEM: {"id": "C72-135", "title": "Extending descent-and-mixing methods to a fixed number of players", "jel_code": "C72", "source_id": "OP-0133"}
% !TeX program = xelatex
\documentclass[11pt,a4paper]{article}
\usepackage[margin=23mm]{geometry}
\usepackage{fontspec}
\setmainfont{Latin Modern Roman}
\usepackage{amsmath,amssymb,mathtools}
\usepackage{xcolor,enumitem,booktabs,longtable,array,xurl,hyperref}
\definecolor{muted}{HTML}{53636C}
\hypersetup{colorlinks=true,urlcolor=blue,linkcolor=blue}
\setlength{\parindent}{0pt}
\setlength{\parskip}{5pt}
\newcommand{\R}{\mathbb R}
\newcommand{\E}{\mathbb E}
\newcommand{\N}{\mathbb N}
\newcommand{\ind}{\mathbf 1}
\newcommand{\Var}{\operatorname{Var}}
\newcommand{\Cov}{\operatorname{Cov}}
\newcommand{\supp}{\operatorname{supp}}
\DeclareMathOperator*{\argmin}{arg\,min}
\DeclareMathOperator*{\argmax}{arg\,max}
\newcommand{\entrymeta}[3]{{\sffamily\small\color{muted}\textbf{\texttt{#1}}\quad|\quad #2\par #3\par}}
\newcommand{\Problemref}[1]{\texttt{#1}}
\newcommand{\JELref}[2]{\texttt{#2} (JEL \texttt{#1})}
\begin{document}
\section*{Extending descent-and-mixing methods to a fixed number of players}
\textbf{Problem ID:} \texttt{C72-135}\quad \textbf{JEL:} \texttt{C72}\par
% BEGIN ECONOMICS STATEMENT
\label{problem:OP-0133}
% GAME_THEORY_PROBLEM: {"local_id": "GT-003", "original_number": 3, "kind": "R", "entry_kind": "statement_only", "published": false, "review_required": true, "category": "game_theory"}
\label{game:GT-003}
{\sffamily\small\color{muted}
\textbf{GT-003} | Algorithmic equilibrium and complexity\\
\textbf{R} | Editorial quantitative version of an explicitly published methodological agenda.
\par}
\subsection*{Definitions and mathematical statement}
Fix an integer $k\geq3$. A game has action sets $S_i=\{1,\ldots,m_i\}$ and explicitly tabulated rational payoffs $u_i:\prod_{j=1}^kS_j\to[0,1]$. Its mixed strategy space is $X=\prod_i\Delta_{m_i}$, and $u_i(x)$ denotes the expectation under independent randomization. Define
\[
 f_G(x)=\max_{1\leq i\leq k}\max_{a_i\in S_i}
 \bigl(u_i(a_i,x_{-i})-u_i(x)\bigr).
\]
The source proposes extending the Tsaknakis--Spirakis (TS) approach, which searches for approximate stationary points of a Nash-residual function and then mixes derived candidates. A concrete target is: for every fixed $k$ and rational $\delta>0$, construct an explicitly described extension of this approach that outputs rational $x\in X$ with
\[
 f_G(x)\leq \frac13+\delta
\]
using a number of bit operations polynomial in the explicit game's length $L$ and $1/\delta$, with the polynomial allowed to depend on $k$.
\subsection*{Short English statement}
Can the descent-and-mixing method give a comparably strong approximation for games with three or more players?
\subsection*{Variants and scope boundaries}
The numerical target $1/3+\delta$ is selected here; the source does not assert it as a canonical multiplayer conjecture. A weaker target may replace $1/3$ by a specified $\alpha_k<1$. A theorem must describe its search and mixing rules: ``trial/search-style'' alone is not a mathematical algorithm class. Complexity is measured in the full payoff table, whose size grows with $\prod_i m_i$.
\subsection*{Source relationship and status}
The original catalog's unspecified ``nontrivial'' and ``comparable'' guarantees are replaced by a declared quantitative benchmark. This is a research-agenda formalization, not an attributed existence conjecture with that exact constant.
\subsection*{References}
\begingroup\footnotesize\raggedright
{\footnotesize\raggedright
\textbf{[S1]} Hanyu Li, Wenhan Huang, Zhijian Duan, David Henry Mguni, Kun Shao, Jun Wang, and Xiaotie Deng (2023). \emph{A survey on algorithms for Nash equilibria in finite normal-form games}. Locator: concluding ``Open problems and future directions,'' theoretical item 2; Section 3.4 for TS methodology. \url{https://arxiv.org/pdf/2312.11063}\par}
\par\endgroup

\subsection*{Common conventions: Source mathematical conventions}

\textbf{Finite games.} For a finite set $A$, $\Delta(A)$ is its probability
simplex. Players choose independent mixed strategies in Nash equilibrium;
correlation is allowed only when explicitly part of the solution concept.
In a game with payoff functions $u_i$ and strategy profile $x$, an additive
$\varepsilon$ Nash equilibrium satisfies
\[
 u_i(x)\geq u_i(x_i',x_{-i})-\varepsilon
 \quad\text{for every player $i$ and every admissible deviation $x_i'$}.
\]
For numerical approximation, payoffs are normalized as locally specified.
An $\varepsilon$ payoff residual is distinct from distance at most
$\varepsilon$ to the set of exact equilibria. Point convergence, convergence
of empirical play, and convergence in distance to an equilibrium set are
also distinct.

\textbf{Computational representations.} Unless stated otherwise, a complexity
question takes rational data with binary encoding; polynomial time is measured
in total bit length. A fixed-accuracy polynomial algorithm is not necessarily
polynomial in $1/\varepsilon$, and neither is necessarily polynomial in
$\log(1/\varepsilon)$. Succinct games and value, demand, or sampling oracles
must declare their interfaces. Exact real solutions require an explicit
representation; an unspecified list of arbitrary real numbers is not a
finite computational output. A claimed algorithmic barrier is meaningful
only for its specified model and reductions.

\textbf{Dynamic games.} Histories, information, observation, and allowable
randomization are part of the model. A stationary strategy depends only on
the current state; a positional strategy in a deterministic graph game
chooses one action at each controlled vertex. Nash equilibrium at an initial
state and equilibrium after every history are different requirements.
An expectation of a pathwise $\liminf$ payoff, a $\liminf$ of expected averages,
a limiting expected average, and a uniform finite-horizon equilibrium are
different criteria. None is silently substituted for another.

\textbf{Allocation and mechanisms.} A complete allocation assigns every item
exactly once unless otherwise stated. Zero-valued goods, negative values,
capacity constraints, feasibility, lotteries, and copies of goods affect
fairness definitions and are specified locally. Dominant-strategy incentive
compatibility, Bayesian incentive compatibility, universal truthfulness,
and truthfulness in expectation impose different inequalities. Whenever
an objective ratio has zero denominator, its convention is stated or those
instances are excluded.

\textbf{Infinite-dimensional models.} Expectations, integrals, stochastic
equations, and optimization objectives require their local regularity,
integrability, filtrations, and admissible domains. A backward--forward
mean-field system is a boundary-value problem; it is not an initial-value
semiflow without an additional construction. Long-time, large-population,
and vanishing-noise limits require an order of limits and a convergence
topology. If a minimum need not be attained, the objective is its infimum
or supremum and the approximation requirement is stated separately.
% END ECONOMICS STATEMENT
\end{document}
